\documentclass[11pt,a4paper,fontset=fandol]{ctexart}

\usepackage[a4paper,top=2.4cm,bottom=2.6cm,left=2.35cm,right=2.35cm,headheight=18pt,headsep=0.55cm,footskip=26pt]{geometry}
\usepackage{amsmath,amssymb}
\usepackage{fancyhdr}
\usepackage{lastpage}
\usepackage{enumitem}
\usepackage{titlesec}
\usepackage{xcolor}
\usepackage{hyperref}
\usepackage{bookmark}
\usepackage[most]{tcolorbox}
\usepackage{setspace}
\usepackage{array}

\hypersetup{
  colorlinks=true,
  linkcolor=blue!50!black,
  urlcolor=blue!50!black
}

\setstretch{1.16}
\setlength{\parindent}{2em}
\setlength{\parskip}{0.32em}
\setlist[enumerate]{itemsep=0.35em, topsep=0.35em, leftmargin=2.3em}
\setlist[itemize]{itemsep=0.25em, topsep=0.25em, leftmargin=2em}

\definecolor{mainblue}{RGB}{36,83,140}
\definecolor{lightblue}{RGB}{241,246,252}
\definecolor{lightgray}{RGB}{246,246,246}
\definecolor{rulegray}{RGB}{110,118,128}
\definecolor{softgreen}{RGB}{236,247,240}
\definecolor{softyellow}{RGB}{254,249,229}

\titleformat{\section}
  {\Large\bfseries\color{mainblue}}
  {\thesection.}
  {0.45em}
  {}
  [\vspace{0.25em}\color{rulegray}\titlerule]
\titlespacing*{\section}{0pt}{1.2em}{0.8em}
\titleformat{\subsection}{\large\bfseries\color{mainblue}}{\thesubsection}{0.5em}{}

\pagestyle{fancy}
\fancyhf{}
\fancyhead[L]{\small 染色方法与棋盘问题专题目录页}
\fancyhead[C]{\small 组合专题资料索引}
\fancyhead[R]{\small 刘文博}
\fancyfoot[C]{\small 第 \thepage 页 / 共 \pageref{LastPage} 页}
\renewcommand{\headrulewidth}{0.55pt}
\renewcommand{\footrulewidth}{0.35pt}

\newtcolorbox{goalbox}[1]{
  breakable,
  colback=lightblue,
  colframe=mainblue,
  boxrule=0.7pt,
  arc=1mm,
  title=\textbf{#1},
  fonttitle=\bfseries
}

\newtcolorbox{infobox}[1]{
  breakable,
  colback=softgreen,
  colframe=green!40!black,
  boxrule=0.65pt,
  arc=1mm,
  title=\textbf{#1},
  fonttitle=\bfseries
}

\newtcolorbox{warnbox}[1]{
  breakable,
  colback=softyellow,
  colframe=orange!60!black,
  boxrule=0.65pt,
  arc=1mm,
  title=\textbf{#1},
  fonttitle=\bfseries
}

\begin{document}

\thispagestyle{empty}
\begin{titlepage}
\centering
\vspace*{0.9cm}

\begin{tcolorbox}[
  enhanced,
  width=0.92\textwidth,
  colback=white,
  colframe=mainblue,
  boxrule=1pt,
  arc=0mm,
  left=6mm,
  right=6mm,
  top=8mm,
  bottom=8mm
]
\centering

{\fontsize{24}{30}\selectfont\bfseries 染色方法与棋盘问题专题目录页\par}
\vspace{0.65cm}
{\fontsize{17}{22}\selectfont 讲义、题单与周测的使用说明\par}

\vspace{1cm}
\begin{tcolorbox}[colback=lightblue,colframe=mainblue,width=0.84\textwidth,boxrule=1pt]
\centering
{\large 适合对象：已经掌握抽屉原理、不变量与基础组合，准备系统学习“用染色做不可能性证明”的学生}\\[0.35em]
{\normalsize 本目录页帮助把染色方法与棋盘问题专题资料串成一个完整训练包}
\end{tcolorbox}

\vspace{1.0cm}
\renewcommand{\arraystretch}{1.42}
\begin{tabular}{>{\bfseries}p{3.5cm}p{8cm}}
专题名称： & 染色方法与棋盘问题 \\
资料组成： & 目前已完成课堂讲义，后续可继续补提高训练题单、周测 A/B \\
使用方式： & 先学讲义，做课堂练习与讲义例题，再逐步扩展到题单与周测 \\
编译方式： & XeLaTeX \\
\end{tabular}

\vspace{0.9cm}
{\large \today\par}
\end{tcolorbox}
\end{titlepage}

\tableofcontents
\newpage

\section{专题包目前包含哪些资料}

\begin{goalbox}{当前阶段的核心资料}
\begin{enumerate}
  \item 课堂讲义：系统讲解黑白染色、多色染色、棋盘覆盖、路径奇偶性与不可能性证明；
  \item 课堂练习与解答：帮助在讲义内部完成第一次巩固。
\end{enumerate}
\end{goalbox}

\begin{infobox}{当前已完成文件}
\begin{itemize}
  \item 讲义：\texttt{coloring\_chessboard\_handout.tex} / \texttt{coloring\_chessboard\_handout.pdf}
\end{itemize}
\end{infobox}

\section{建议学习顺序}

\begin{goalbox}{推荐顺序}
\begin{enumerate}
  \item 先完整学习讲义，重点吃透“为什么要染色”“每一种棋子对颜色有什么固定作用”；
  \item 再独立完成讲义中的课堂练习，并对照解答订正；
  \item 等后续题单与周测补齐后，再进入系统检测阶段。
\end{enumerate}
\end{goalbox}

\begin{warnbox}{不建议的做法}
\begin{itemize}
  \item 一看到棋盘题就先试着硬拼，而不先想颜色分布；
  \item 染了颜色，却没有说明“每一步为什么一定这样变化”；
  \item 把“黑白格数相等”误当成充分条件；
  \item 只会黑白染色，不会根据题目改用三色或更多颜色。
\end{itemize}
\end{warnbox}

\section{这份讲义主要解决什么问题}

\begin{goalbox}{学完讲义应重点掌握的内容}
\begin{enumerate}
  \item 会用黑白染色处理最基本的骨牌覆盖题；
  \item 会用路径换色规律判断某些走法是否可能；
  \item 会根据覆盖块形状设计三色染色；
  \item 会区分“必要条件”和“充分条件”。
\end{enumerate}
\end{goalbox}

\section{一个推荐的 3 次训练安排}

\begin{goalbox}{3 次完成讲义阶段训练}
\begin{enumerate}
  \item 第 1 次：学习讲义前半部分，重点理解黑白染色与残缺棋盘；
  \item 第 2 次：学习讲义后半部分，重点理解多色染色与路径问题；
  \item 第 3 次：回做讲义练习，重点订正“颜色为什么有效”和“必要条件是否被误用”。
\end{enumerate}
\end{goalbox}

\begin{infobox}{如果孩子基础较好，可以这样压缩}
\begin{itemize}
  \item 第 1 天：讲义 + 前半部分练习；
  \item 第 2 天：讲义后半部分 + 讲义全部练习订正；
  \item 第 3 天：口头复述本专题的核心方法，并回顾与抽屉原理、不变量的联系。
\end{itemize}
\end{infobox}

\section{专题学习后的自检清单}

\begin{goalbox}{如果下面 7 条都能做到，说明讲义阶段基本过关}
\begin{enumerate}
  \item 我会判断一道棋盘题是否值得先染色；
  \item 我会解释一个 $1\times 2$ 骨牌为什么总覆盖一黑一白；
  \item 我会用黑白格数量判断某些覆盖题不可能；
  \item 我会用颜色交替判断某些路径题能不能到达；
  \item 我会理解三色染色为什么适合 $1\times 3$ 长条；
  \item 我知道“颜色数相等”常常只是必要条件；
  \item 我会在解答后检查：这种染色到底利用了什么固定性质。
\end{enumerate}
\end{goalbox}

\section{给家长的使用建议}

\begin{warnbox}{更有效的带练方式}
\begin{itemize}
  \item 不要一开始就给结论，先让孩子自己说“为什么想染色”；
  \item 如果卡住，可以只提示“这块骨牌每次会碰到哪些颜色”；
  \item 订正时重点看孩子有没有把“固定覆盖关系”说清楚；
  \item 做完后请孩子口头复述：这题是用黑白染色、三色染色，还是颜色交替解决的。
\end{itemize}
\end{warnbox}

\begin{infobox}{和前后专题的衔接}
这个专题和\textbf{抽屉原理与极值思想}、\textbf{数论方法中的组合构造}联系都很紧。前者提供“分类和存在性证明”的视角，后者提供“模分类和不变量”的视角，而染色方法可以看作它们在棋盘题中的自然延伸。
\end{infobox}

\end{document}
